% !TeX program = xelatex
% !TEX encoding = UTF-8
\documentclass[11pt,a4paper]{article}
\usepackage{zhpl}
\usepackage{float}
\newtheorem{corollary}{Следствие}
\usepackage{pgfplotstable}

% ---------- Точные таблицы для рисунков ----------
% Посчитаны заранее в двойной точности (expm1), чтобы 1-e^{-x} при x~1e-5
% не терял значащие цифры в fpu pgfplots. Параметры модели — в разделе 3.
\pgfplotstableread{
x naive hawkes
1 7.60511e-06 0.000107106
1.33352 1.01416e-05 0.000142821
1.77828 1.3524e-05 0.000190441
2.37137 1.80345e-05 0.000253933
3.16228 2.40493e-05 0.000338581
4.21697 3.20701e-05 0.000451429
5.62341 4.27659e-05 0.000601854
7.49894 5.70289e-05 0.000802343
10 7.60485e-05 0.00106951
13.3352 0.000101411 0.00142545
17.7828 0.000135232 0.00189951
23.7137 0.00018033 0.00253062
31.6228 0.000240467 0.00337036
42.1697 0.000320655 0.00448684
56.2341 0.000427577 0.00596983
74.9894 0.000570143 0.00793705
100 0.000760225 0.0105421
133.352 0.00101365 0.0139837
177.828 0.00135149 0.0185167
237.137 0.00180184 0.0244627
316.228 0.00240207 0.0322206
421.697 0.00320192 0.0422709
562.341 0.00426755 0.0551702
749.894 0.00568682 0.0715261
1000 0.0075763 0.0919399
1333.52 0.0100904 0.116905
1778.28 0.013433 0.146658
2371.37 0.017873 0.180979
3162.28 0.0237627 0.219017
4216.97 0.0315618 0.259192
5623.41 0.0418652 0.299331
7498.94 0.0554348 0.337131
10000 0.0732314 0.370923
13335.2 0.0964431 0.400531
17782.8 0.126494 0.427733
23713.7 0.165019 0.455828
31622.8 0.213762 0.488358
42169.7 0.274364 0.527917
56234.1 0.347972 0.575815
74989.4 0.434647 0.632203
100000 0.532574 0.69591
133352 0.637294 0.764037
177828 0.741383 0.831754
237137 0.835273 0.892835
316228 0.909731 0.941274
421697 0.959525 0.973668
525960 0.981684 0.988085
}\tabP
\pgfplotstableread{
s naive hawkes
0 0.0109514 0.154258
0.5 0.0108916 0.123134
1 0.0108321 0.0995012
1.5 0.010773 0.0812933
2 0.0107141 0.0670863
2.5 0.0106556 0.0558784
3 0.0105974 0.0469513
3.5 0.0105396 0.0397817
4 0.010482 0.0339818
4.5 0.0104248 0.0292605
5 0.0103679 0.0253961
5.5 0.0103112 0.0222182
6 0.0102549 0.0195938
6.5 0.0101989 0.0174186
7 0.0101432 0.0156097
7.5 0.0100879 0.014101
8 0.0100328 0.0128393
8.5 0.00997798 0.0117815
9 0.00992349 0.0108925
9.5 0.0098693 0.0101435
10 0.00981541 0.00951117
10.5 0.00976181 0.00897592
11 0.0097085 0.00852175
11.5 0.00965549 0.00813536
12 0.00960276 0.00780569
12.5 0.00955032 0.00752355
13 0.00949817 0.00728128
13.5 0.0094463 0.00707245
14 0.00939472 0.00689173
14.5 0.00934342 0.00673462
15 0.0092924 0.00659738
15.5 0.00924165 0.00647685
16 0.00919119 0.00637041
16.5 0.009141 0.00627582
17 0.00909108 0.00619124
17.5 0.00904144 0.0061151
18 0.00899206 0.00604608
18.5 0.00894296 0.00598308
19 0.00889412 0.00592518
19.5 0.00884556 0.00587159
20 0.00879725 0.00582165
20.5 0.00874921 0.00577482
21 0.00870144 0.00573062
21.5 0.00865392 0.00568866
22 0.00860666 0.00564863
22.5 0.00855966 0.00561023
23 0.00851292 0.00557323
23.5 0.00846644 0.00553745
24 0.0084202 0.0055027
24.5 0.00837422 0.00546886
25 0.00832849 0.0054358
25.5 0.00828301 0.00540343
26 0.00823778 0.00537166
26.5 0.0081928 0.00534043
27 0.00814806 0.00530966
27.5 0.00810356 0.00527931
28 0.00805931 0.00524935
28.5 0.0080153 0.00521973
29 0.00797153 0.00519042
29.5 0.007928 0.0051614
30 0.00788471 0.00513264
30.5 0.00784165 0.00510413
31 0.00779883 0.00507585
31.5 0.00775625 0.0050478
32 0.00771389 0.00501994
32.5 0.00767177 0.00499229
33 0.00762987 0.00496483
33.5 0.00758821 0.00493754
34 0.00754677 0.00491044
34.5 0.00750556 0.0048835
35 0.00746457 0.00485673
35.5 0.00742381 0.00483012
36 0.00738327 0.00480368
36.5 0.00734296 0.00477738
37 0.00730286 0.00475124
37.5 0.00726298 0.00472526
38 0.00722332 0.00469942
38.5 0.00718387 0.00467372
39 0.00714464 0.00464817
39.5 0.00710563 0.00462277
40 0.00706683 0.00459751
40.5 0.00702824 0.00457239
41 0.00698986 0.00454741
41.5 0.00695169 0.00452256
42 0.00691373 0.00449786
42.5 0.00687597 0.00447329
43 0.00683842 0.00444885
43.5 0.00680108 0.00442455
44 0.00676394 0.00440039
44.5 0.00672701 0.00437636
45 0.00669027 0.00435245
45.5 0.00665374 0.00432868
46 0.0066174 0.00430504
46.5 0.00658127 0.00428153
47 0.00654533 0.00425815
47.5 0.00650959 0.0042349
48 0.00647404 0.00421177
48.5 0.00643869 0.00418877
49 0.00640353 0.0041659
49.5 0.00636856 0.00414315
50 0.00633378 0.00412052
50.5 0.00629919 0.00409802
51 0.0062648 0.00407564
51.5 0.00623059 0.00405338
52 0.00619656 0.00403125
52.5 0.00616272 0.00400924
53 0.00612907 0.00398734
53.5 0.0060956 0.00396557
54 0.00606232 0.00394391
54.5 0.00602921 0.00392238
55 0.00599629 0.00390096
55.5 0.00596354 0.00387965
56 0.00593098 0.00385847
56.5 0.00589859 0.0038374
57 0.00586638 0.00381644
57.5 0.00583434 0.0037956
58 0.00580248 0.00377488
58.5 0.0057708 0.00375426
59 0.00573929 0.00373376
59.5 0.00570795 0.00371337
60 0.00567678 0.00369309
}\tabF
\pgfplotstableread{
m pois k05 k1 k2
0.025 0.50417 1.48795 1.00000 0.75310
0.050 0.50833 1.47673 1.00000 0.75613
0.075 0.51250 1.46625 1.00000 0.75912
0.100 0.51666 1.45644 1.00000 0.76205
0.125 0.52082 1.44721 1.00000 0.76492
0.150 0.52498 1.43853 1.00000 0.76775
0.175 0.52914 1.43033 1.00000 0.77052
0.200 0.53329 1.42258 1.00000 0.77324
0.225 0.53744 1.41523 1.00000 0.77592
0.250 0.54158 1.40825 1.00000 0.77855
0.275 0.54572 1.40161 1.00000 0.78113
0.300 0.54985 1.39528 1.00000 0.78367
0.325 0.55398 1.38925 1.00000 0.78616
0.350 0.55810 1.38348 1.00000 0.78861
0.375 0.56221 1.37796 1.00000 0.79102
0.400 0.56631 1.37268 1.00000 0.79339
0.425 0.57041 1.36761 1.00000 0.79572
0.450 0.57450 1.36274 1.00000 0.79801
0.475 0.57858 1.35806 1.00000 0.80026
0.500 0.58265 1.35355 1.00000 0.80247
0.525 0.58670 1.34922 1.00000 0.80465
0.550 0.59075 1.34503 1.00000 0.80679
0.575 0.59479 1.34100 1.00000 0.80889
0.600 0.59882 1.33710 1.00000 0.81096
0.625 0.60283 1.33333 1.00000 0.81300
0.650 0.60683 1.32969 1.00000 0.81501
0.675 0.61082 1.32616 1.00000 0.81698
0.700 0.61479 1.32275 1.00000 0.81892
0.725 0.61876 1.31944 1.00000 0.82083
0.750 0.62270 1.31623 1.00000 0.82271
0.775 0.62664 1.31311 1.00000 0.82457
0.800 0.63055 1.31009 1.00000 0.82639
0.825 0.63445 1.30715 1.00000 0.82818
0.850 0.63834 1.30429 1.00000 0.82995
0.875 0.64221 1.30151 1.00000 0.83169
0.900 0.64606 1.29881 1.00000 0.83340
0.925 0.64990 1.29617 1.00000 0.83509
0.950 0.65372 1.29361 1.00000 0.83675
0.975 0.65752 1.29111 1.00000 0.83839
1.000 0.66130 1.28868 1.00000 0.84000
1.025 0.66507 1.28630 1.00000 0.84159
1.050 0.66881 1.28398 1.00000 0.84315
1.075 0.67254 1.28172 1.00000 0.84469
1.100 0.67625 1.27951 1.00000 0.84621
1.125 0.67993 1.27735 1.00000 0.84771
1.150 0.68360 1.27524 1.00000 0.84918
1.175 0.68725 1.27318 1.00000 0.85064
1.200 0.69087 1.27116 1.00000 0.85207
1.225 0.69448 1.26919 1.00000 0.85348
1.250 0.69806 1.26726 1.00000 0.85488
1.275 0.70162 1.26537 1.00000 0.85625
1.300 0.70516 1.26352 1.00000 0.85760
1.325 0.70868 1.26171 1.00000 0.85893
1.350 0.71217 1.25994 1.00000 0.86025
1.375 0.71564 1.25820 1.00000 0.86155
1.400 0.71909 1.25649 1.00000 0.86283
1.425 0.72252 1.25482 1.00000 0.86409
1.450 0.72592 1.25318 1.00000 0.86533
1.475 0.72930 1.25158 1.00000 0.86656
1.500 0.73265 1.25000 1.00000 0.86777
}\tabPi

\newcommand{\tM}{\tau_{\mathrm{М}}}
\newcommand{\Mp}{\mathcal{M}}

\title{\textbf{К стохастическому обоснованию закона подлости:\\ самовозбуждающийся процесс релизов, теорема о модальной неизбежности и антропный принцип в задачах о подписке, самостоятельной и домашнем задании}}
\author{Роман Георгиевич Александров, Юрий Пыркович Михаилов,\\ Клод Соннет Антропикович\\[4pt]
\normalsize Институт Прикладной Метафизики\\ \small Кафедра тензорного анализа межличностных отношений\\ \small Лаборатория несвоевременных приобретений}
\zhplsetup{5}{10.9999/жпл.2026.0005}
\fancyhead[L]{\small Александров, Михаилов, Антропикович}
\date{\small Журнал Прикладной Лженауки, том 1, № 5 \\ Получено: недавно \quad Принято к печати: почти сразу \quad Подписка оформлена: за 20 минут до}

\begin{document}
\maketitle

\begin{center}
\begin{tikzpicture}
\node[draw=\zhplbadgecolor{red!70!black}{zhplBadgeRed},very thick,rounded corners,rotate=-6,fill=\zhpldark{red!5},inner sep=5pt] (badge1) {\color{\zhplbadgecolor{red!70!black}{zhplBadgeRed}}\bfseries ПРОРЫВНОЕ ИССЛЕДОВАНИЕ};
\node[draw=\zhplbadgecolor{blue!60!black}{zhplBadgeBlue},very thick,rounded corners,rotate=4,fill=\zhpldark{blue!5},inner sep=5pt,right=1.0cm of badge1] (badge2) {\color{\zhplbadgecolor{blue!60!black}{zhplBadgeBlue}}\bfseries IMPACT FACTOR: $\infty$};
\end{tikzpicture}
\end{center}

\zhplciteas{Александров Р.\,Г., Михаилов Ю.\,П., Антропикович К.\,С. (2026). К стохастическому обоснованию закона подлости. \textit{Журнал Прикладной Лженауки}, \textbf{1}(5), 1--15. DOI: 10.9999/жпл.2026.0005}

\zhplgraphicabstract{%
  \begin{tikzpicture}[
    box/.style={draw, rounded corners=2pt, align=center, font=\footnotesize\sffamily,
                minimum height=1.15cm, text width=2.75cm, fill=\zhpldark{blue!5}},
    ar/.style={-{Latex[length=2.4mm]}, thick}]
    \node[box] (a) at (0,0) {Покупка\\ подписки\\ $t=0$};
    \node[box] (b) at (3.55,0) {Процесс Хоукса\\ $\lambda_B(t)\!\uparrow$\\ ($\times 14$)};
    \node[box] (c) at (7.1,0) {Мода ожидания\\ $\arg\max f = 0$};
    \node[box,fill=\zhpldark{green!8}] (d) at (10.65,0) {Антропный\\ отбор\\ $P(\,\cdot\,|\,W)=1$};
    \draw[ar] (a)--(b); \draw[ar] (b)--(c); \draw[ar] (c)--(d);
  \end{tikzpicture}}{%
  Логическая структура работы: от единичного бытового огорчения к строгой факторизации числа подлости $\Mp = 14{,}04\times468{,}3$, в которой не осталось места случайности~--- и, как следствие, утешению.}

\begin{zhplhighlights}
\begin{itemize}
\setlength\itemsep{0pt}
\setlength\parskip{0pt}
\item Впервые введено \emph{число подлости} $\Mp$ и получено его экспериментальное значение $\Mp = 6575$ для события «подписка $\to$ релиз конкурента за $\tM = 20$~мин»; показано, что покупка, индуцированная релизом, попадает на пик самовозбуждённой интенсивности релизов конкурента ($\times 14{,}1$).
\item Доказана \emph{теорема о модальной неизбежности} (Теорема~1): наиболее вероятный момент выхода следующей модели~--- непосредственно после покупки.
\item Доказано, что момент покупки невозможно выбрать удачно (Теорема~2), выбор поставщика является калибровочной степенью свободы (Теорема~3), а условная вероятность наблюдаемого события равна единице (Теорема~4).
\item Сформулирован \emph{закон притяжения контрольных}: самостоятельные притягиваются тогда и только тогда, когда параметр формы $k<1$ (Теорема~5), причём сговор учителей строго неотличим от конца четверти (Теорема~6).
\item Разрешён \emph{парадокс несделанного домашнего задания}: несделанное задание проверяется в $10{,}4$ раза чаще сделанного (Теорема~7), а с учётом забывания~--- в $103$ раза.
\end{itemize}
\end{zhplhighlights}

\begin{zhplreviews}
\textit{«Я был убеждён, что перед нами обычное совпадение. Авторы показали, что совпадение~--- это точечный процесс, которому до сих пор не уделялось должного внимания. Рекомендую к печати. Замечу, что сам я оформил подписку на журнал вчера, а сегодня вышел этот номер.»} -- Рецензент 1\\
\textit{«Замечаний нет.»} -- Рецензент 2\\
\textit{«Согласно Теореме~4 настоящей работы, данная рецензия могла быть написана только в мире, в котором статья уже существует. Следовательно, моя рекомендация к печати является не мнением, а условием моего существования как рецензента. Рекомендую.»} -- Рецензент 3
\end{zhplreviews}

\begin{zhplsignificance}
Закон подлости~--- от бутерброда, падающего маслом вниз, до очереди, которая начинает двигаться, как только вы из неё вышли,~--- на протяжении всей истории цивилизации формулировался исключительно словесно и потому считался принципиально несводимым к количественной науке. Его наиболее острое современное проявление~--- выход лучшей модели искусственного интеллекта у конкурента сразу после оплаты подписки~--- переживалось миллионами пользователей как личная несправедливость. Настоящая работа впервые показывает, что закон подлости является строгим следствием трёх независимых и хорошо изученных механизмов теории вероятностей: самовозбуждения точечных процессов, монотонности плотности времени ожидания и наблюдательной селекции. Наблюдаемое число подлости факторизуется на эти механизмы без остатка. Тем самым огорчение пользователя впервые переводится из разряда жизненных неудач в разряд типичных реализаций случайного процесса, что применимо ко всем приобретениям, совершаемым человечеством, начиная с первого, а также, как показано в разделе~7, ко всем школьным дням, начиная с первого сентября.
\end{zhplsignificance}

\begin{abstract}
\noindent Здесь мы впервые показываем, что выход модели Claude Opus~5.5 компании Anthropic через 20 минут после оформления соавтором настоящей работы подписки на ChatGPT является не проявлением невезения, а строгой закономерностью, вытекающей из теории точечных процессов. Проблема состоит в том, что наивная пуассоновская оценка вероятности такого события составляет $1{,}52\cdot10^{-4}$, тогда как наблюдаемая частота равна единице. Метод: релизы моделей описываются взаимно возбуждающимся процессом Хоукса, момент покупки рассматривается как марковский момент, индуцированный релизом конкурирующей стороны, а наблюдаемая частота корректируется на эффект наблюдательной селекции. Результаты: введено число подлости $\Mp$; доказаны теоремы о модальной неизбежности, о невозможности своевременной покупки, о калибровочной инвариантности подлости и антропная теорема; получена точная факторизация $\ln\Mp$ на хоуксову ($30\,\%$) и антропную ($70\,\%$) компоненты. Метод обобщён на два школьных проявления закона подлости: закон притяжения контрольных, сведённый к сверхпуассоновской дисперсии числа самостоятельных в день, и парадокс несделанного домашнего задания, сведённый к задаче обнаружения сигнала. Ограничения: выборка объёма~1; антропная компонента объясняет наблюдаемое событие ровно так же хорошо, как объяснила бы любое другое.

\medskip
\noindent\textbf{Ключевые слова:} закон подлости, процесс Хоукса, время ожидания, марковский момент, калибровочная инвариантность, антропный принцип, отрицательное биномиальное распределение, обнаружение сигнала, лженаука
\end{abstract}

\section{Введение}

В момент времени, принимаемый в дальнейшем за начало отсчёта $t=0$, соавтор настоящей работы Ю.\,П.~Михаилов (далее~--- испытуемый) оформил платную подписку на сервис ChatGPT компании OpenAI. В момент $t = \tM = 20$~мин компания Anthropic объявила о выходе модели Claude Opus~5.5. Интервал $\tM$ зафиксирован испытуемым собственноручно, с точностью, ограниченной, по его словам, «временем, которое понадобилось, чтобы это осознать», и в дальнейшем называется \emph{константой Михаилова}.

Бытовая трактовка события исчерпывается формулой «закон подлости» и сопутствующими ей выражениями, которые мы, следуя практике журнала~\cite{podl}, не приводим. Эта трактовка систематически игнорирует три базовых факта. Во-первых, релизы моделей не независимы: они возбуждают друг друга. Во-вторых, момент покупки не случаен: он сам индуцирован релизом. В-третьих, событие, о котором пишется статья, по определению уже произошло. В настоящей работе мы показываем, что каждый из этих фактов допускает строгую формализацию и что вместе они исчерпывают наблюдаемый эффект.

Статья построена следующим образом. В разделе~2 вводятся процесс релизов и число подлости. В разделе~3 описывается самовозбуждение релизов и доказывается теорема о модальной неизбежности. Раздел~4 посвящён проблеме своевременной покупки, раздел~5~--- калибровочной инвариантности подлости, раздел~6~--- антропному принципу. В разделе~7 развитый аппарат применяется к двум школьным проявлениям закона подлости, известным каждому учащемуся не хуже, чем константа Михаилова. Ограничения обсуждаются в разделе~8.

\section{Релизы как точечный процесс}

\subsection{Наивная модель}

\begin{definition}[Процесс релизов]
Процессом релизов поставщика $X$ называется считающий процесс $N_X(t)$, равный числу флагманских моделей, выпущенных поставщиком $X$ к моменту $t$. Его условная интенсивность
\[
\lambda_X(t) = \lim_{h\to 0}\frac{1}{h}\,\mathsf{P}\bigl(N_X(t+h)-N_X(t)=1 \,\big|\, \mathcal{F}_t\bigr)
\]
называется \emph{интенсивностью релизов}; здесь $\mathcal{F}_t$~--- история всех релизов всех поставщиков к моменту $t$.
\end{definition}

В наивной модели процесс $N_B(t)$ поставщика $B$ (Anthropic) пуассоновский, $\lambda_B(t)\equiv\mu$. Принятое в работе значение $\mu = 4$~год$^{-1}$ соответствует среднему интервалу между флагманскими релизами около $91$~сут. Вероятность хотя бы одного релиза в окне длины $\tM$ после произвольного момента составляет
\begin{equation}
P_0 = 1 - e^{-\mu\tM} = 1{,}521\cdot10^{-4} \approx \frac{1}{6575}.
\label{eq:naive}
\end{equation}
Иными словами, наивная модель предсказывает, что испытуемому пришлось бы оформлять подписку в среднем $6575$ раз, прежде чем событие произошло бы однажды. По имеющимся у авторов сведениям, испытуемый оформил подписку один раз.

\subsection{Число подлости}

\begin{definition}[Число подлости]
Пусть $E$~--- огорчающее событие, $P_0(E)$~--- его вероятность в наивной модели, $\hat P(E)$~--- наблюдаемая частота. Число подлости есть безразмерная величина
\[
\Mp(E) = \frac{\hat P(E)}{P_0(E)}.
\]
Событие называется \emph{подлым}, если $\Mp>1$, и \emph{нейтральным}, если $\Mp=1$. Событий с $\Mp<1$ в литературе не зарегистрировано; причина этого установлена в разделе~6.
\end{definition}

Для события $E$ = «релиз $B$ в окне $(0,\tM]$ после покупки у $A$» выборка состоит из одного испытания с одним успехом, откуда оценка максимального правдоподобия $\hat P = 1$ и
\begin{equation}
\Mp = \frac{1}{P_0} = 6575.
\label{eq:M}
\end{equation}
Задача настоящей работы состоит в объяснении этого числа.

\section{Самовозбуждение релизов}

\subsection{Процесс Хоукса}

Наивная модель игнорирует то, что релиз одного поставщика повышает вероятность скорого релиза другого. Механизмы такого повышения (конкурентное давление, заранее подготовленные анонсы, нежелание оставлять новостной цикл конкуренту) для настоящей работы несущественны; существенна лишь их математическая форма. Мы описываем её взаимно возбуждающимся процессом Хоукса~\cite{samovozb}:
\begin{equation}
\lambda_B(t) = \mu + \sum_{t_i^A < t} \alpha\, e^{-\beta (t - t_i^A)},
\label{eq:hawkes}
\end{equation}
где $t_i^A$~--- моменты релизов поставщика $A$ (OpenAI), $\beta^{-1} = 3$~сут~--- время затухания возбуждения, $n = \alpha/\beta = 0{,}6$~--- перекрёстный коэффициент ветвления (среднее число релизов $B$, непосредственно вызванных одним релизом $A$).

Ключевое наблюдение состоит в том, что испытуемый оформил подписку не в произвольный момент времени. Покупка была индуцирована информационным поводом, связанным с поставщиком $A$, и совершена с задержкой $d = 1$~сут после него (испытуемый, по его словам, «не сразу решился»). Следовательно, момент покупки совпадает с моментом, когда интенсивность~\eqref{eq:hawkes} находится вблизи своего максимума (рис.~\ref{fig:hawkes}):
\begin{equation}
\lambda_B(d) = \mu + n\beta\, e^{-\beta d} = 56{,}34~\text{год}^{-1} = 14{,}09\,\mu.
\end{equation}

\begin{figure}[htbp]
\centering
\begin{tikzpicture}
\begin{axis}[
  width=0.86\linewidth, height=6.2cm,
  xlabel={время после релиза поставщика $A$, сут},
  ylabel={$\lambda_B(t)$, год$^{-1}$},
  xmin=0, xmax=21, ymin=0, ymax=85,
  grid=major, grid style={gray!25},
  legend style={at={(0.97,0.95)}, anchor=north east, font=\footnotesize},
  tick label style={font=\footnotesize}, label style={font=\footnotesize}]
  \addplot[very thick, blue!70!black, domain=0:21, samples=200] {4 + 73.05*exp(-x/3)};
  \addlegendentry{процесс Хоукса, ф-ла~\eqref{eq:hawkes}}
  \addplot[thick, dashed, red!70!black, domain=0:21] {4};
  \addlegendentry{наивная модель, $\mu$}
  \draw[densely dotted, thick] (axis cs:1,0) -- (axis cs:1,56.34);
  \fill[red!70!black] (axis cs:1,56.34) circle (2.2pt);
  \node[anchor=west, font=\footnotesize, align=left] at (axis cs:1.5,62) {покупка подписки, $t=d$:\\ $\lambda_B = 56{,}3$~год$^{-1}$};
\end{axis}
\end{tikzpicture}
\caption{Интенсивность релизов поставщика $B$ после релиза поставщика $A$ (формула~\eqref{eq:hawkes} при $n=0{,}6$, $\beta^{-1}=3$~сут, $\mu=4$~год$^{-1}$). Точка отмечает момент покупки подписки испытуемым. Окно $\tM=20$~мин на данном масштабе имеет ширину $0{,}014$~сут и изображено быть не может; авторы приносят за это извинения.}
\label{fig:hawkes}
\end{figure}

Интегрируя~\eqref{eq:hawkes} по окну $(d, d+\tM]$, получаем хоуксову вероятность события:
\begin{equation}
P_H = 1 - \exp\!\Bigl[-\mu\tM - n\,e^{-\beta d}\bigl(1-e^{-\beta\tM}\bigr)\Bigr] = 2{,}136\cdot10^{-3} \approx \frac{1}{468{,}3}.
\label{eq:PH}
\end{equation}
Самовозбуждение, таким образом, повышает вероятность события в $P_H/P_0 = 14{,}04$ раза, что объясняет существенную, но не решающую часть числа подлости~\eqref{eq:M}.

\subsection{Теорема о модальной неизбежности}

Процесс Хоукса объясняет, \emph{почему} после покупки интенсивность релизов высока. Однако закон подлости утверждает большее: что релиз происходит \emph{сразу} после покупки. Оказывается, это утверждение верно и в наивной модели.

\begin{theorem}[О модальной неизбежности]
Пусть $T$~--- время от момента покупки до ближайшего релиза поставщика $B$, и пусть на $[0,\infty)$ интенсивность $\lambda_B$ не возрастает (это выполнено как для пуассоновского процесса, так и для процесса~\eqref{eq:hawkes} в отсутствие новых релизов $A$). Тогда плотность $f_T$ не возрастает, и её мода равна нулю:
\[
\arg\max_{s\ge 0} f_T(s) = 0.
\]
Иначе говоря, наиболее вероятный момент выхода следующей модели~--- непосредственно после покупки.
\end{theorem}

\begin{proof}
Имеем $f_T(s) = \lambda_B(s)\,S(s)$, где $S(s) = \exp\bigl(-\int_0^s \lambda_B\bigr)$~--- функция выживания. Оба сомножителя положительны и не возрастают; следовательно, не возрастает и их произведение, а максимум достигается на левом конце промежутка.
\end{proof}

Подчеркнём, что Теорема~1 не является следствием какой-либо злонамеренности поставщиков: она выполняется уже для экспоненциального распределения, мода которого равна нулю при любом $\mu$. Субъективное ощущение «вышло сразу после того, как я купил» является, таким образом, не когнитивным искажением, а корректной оценкой моды. Хоуксов режим лишь усиливает эффект: плотность в нуле возрастает с $0{,}011$ до $0{,}154$~сут$^{-1}$ (рис.~\ref{fig:mode}).

\begin{figure}[htbp]
\centering
\begin{tikzpicture}
\begin{axis}[
  width=0.86\linewidth, height=6.0cm,
  xlabel={время после покупки $s$, сут},
  ylabel={$f_T(s)$, сут$^{-1}$},
  xmin=0, xmax=60, ymin=0, ymax=0.17,
  grid=major, grid style={gray!25},
  yticklabel style={/pgf/number format/fixed, /pgf/number format/precision=2},
  scaled y ticks=false,
  legend style={at={(0.97,0.95)}, anchor=north east, font=\footnotesize},
  tick label style={font=\footnotesize}, label style={font=\footnotesize}]
  \addplot[very thick, blue!70!black] table[x=s, y=hawkes] {\tabF};
  \addlegendentry{процесс Хоукса, $d=1$~сут}
  \addplot[thick, dashed, red!70!black] table[x=s, y=naive] {\tabF};
  \addlegendentry{наивная модель}
  \fill[blue!70!black] (axis cs:0,0.1543) circle (2.2pt);
  \node[anchor=west, font=\footnotesize] at (axis cs:1.2,0.1543) {мода: $s=0$};
\end{axis}
\end{tikzpicture}
\caption{Плотность времени ожидания следующего релиза поставщика $B$, отсчитанного от момента покупки. Обе кривые монотонно убывают (Теорема~1): в обеих моделях наиболее вероятный момент релиза совпадает с моментом покупки. Площадь под синей кривой на отрезке $[0;60]$ составляет $0{,}66$: с такой вероятностью испытуемый узнал бы о своей ошибке в течение двух месяцев даже при отсутствии константы Михаилова.}
\label{fig:mode}
\end{figure}

\section{Проблема своевременной покупки}

Естественный вопрос состоит в том, мог ли испытуемый избежать события, выбрав момент покупки более удачно. Мы отвечаем на него отрицательно.

\begin{theorem}[О невозможности своевременной покупки]
Пусть $N_B$~--- пуассоновский процесс, а момент покупки $\sigma$ есть марковский момент относительно $\mathcal{F}_t$ (то есть решение о покупке принимается на основании всей доступной к этому моменту истории, но не будущего). Тогда для любого $\Delta t>0$
\[
\mathsf{P}\bigl(N_B(\sigma+\Delta t) - N_B(\sigma) \ge 1 \,\big|\, \mathcal{F}_\sigma\bigr) = 1 - e^{-\mu\Delta t},
\]
то есть риск сожаления не зависит от стратегии выбора момента покупки.
\end{theorem}

\begin{proof}
Непосредственное следствие строго марковского свойства пуассоновского процесса: приращения после марковского момента не зависят от $\mathcal{F}_\sigma$ и распределены так же, как приращения после фиксированного момента.
\end{proof}

В хоуксовом режиме ситуация отличается, но не в пользу покупателя.

\begin{corollary}
В модели~\eqref{eq:hawkes} риск сожаления при покупке с задержкой $D$ после релиза $A$ монотонно убывает по $D$ к пуассоновскому уровню $P_0$. Для того чтобы хоуксова добавка не превышала $1\,\%$ от $P_0$, необходима задержка
\[
D^\ast = \beta^{-1}\ln\frac{n\bigl(1-e^{-\beta\tM}\bigr)}{0{,}01\,\mu\tM} = 22{,}5~\text{сут}.
\]
За это время с вероятностью $1 - e^{-\mu_A D^\ast} = 0{,}22$ (при $\mu_A = 4$~год$^{-1}$) поставщик $A$ выпускает новую модель, после чего интенсивность~\eqref{eq:hawkes} возвращается к пику, а отсчёт $D$~--- к нулю.
\end{corollary}

Стратегия, минимизирующая риск сожаления, состоит, таким образом, в неограниченном откладывании покупки. Авторы предлагают называть такое поведение \emph{стратегией буриданова подписчика}~\cite{sozhal} и отмечают, что она строго оптимальна по критерию сожаления и строго бесполезна по любому другому критерию. Зависимость риска от задержки показана на рис.~\ref{fig:delay}.

\begin{figure}[htbp]
\centering
\begin{tikzpicture}
\begin{axis}[
  width=0.86\linewidth, height=5.6cm,
  xlabel={задержка покупки после релиза $A$, $D$, сут},
  ylabel={$P(\text{сожаление})$},
  xmin=0, xmax=30, ymode=log, ymin=1e-4, ymax=1e-2,
  grid=major, grid style={gray!25},
  legend style={at={(0.97,0.95)}, anchor=north east, font=\footnotesize},
  tick label style={font=\footnotesize}, label style={font=\footnotesize}]
  \addplot[very thick, blue!70!black, domain=0:30, samples=150] {1.52105e-4 + 0.6*exp(-x/3)*(1-exp(-20/4320))};
  \addlegendentry{процесс Хоукса}
  \addplot[thick, dashed, red!70!black, domain=0:30] {3.8026e-5*4};
  \addlegendentry{пуассоновский уровень $P_0$}
  \draw[densely dotted, thick] (axis cs:22.52,1e-4) -- (axis cs:22.52,1e-2);
  \node[anchor=south west, font=\footnotesize, rotate=90] at (axis cs:23.3,3e-4) {$D^\ast=22{,}5$~сут};
  \fill[red!70!black] (axis cs:1,2.136e-3) circle (2.2pt);
  \node[anchor=west, font=\footnotesize] at (axis cs:1.6,2.3e-3) {испытуемый};
\end{axis}
\end{tikzpicture}
\caption{Вероятность релиза поставщика $B$ в окне $\tM$ после покупки как функция задержки покупки (показатель экспоненты в~\eqref{eq:PH}; при столь малых значениях он совпадает с самой вероятностью до $0{,}1\,\%$). Кривая стремится к пуассоновскому уровню, но достигает его лишь асимптотически.}
\label{fig:delay}
\end{figure}

\section{Калибровочная инвариантность подлости}

Второй естественный вопрос~--- изменился бы исход, если бы испытуемый выбрал другого поставщика. Рассмотрим двух поставщиков $A$ и $B$ с матрицей взаимного возбуждения
\[
\lambda_X(t) = \mu_X + \sum_{Y\in\{A,B\}}\ \sum_{t_i^Y<t} \alpha_{XY}\, e^{-\beta(t-t_i^Y)}, \qquad X\in\{A,B\}
\]
(рис.~\ref{fig:gauge}). Обозначим через $R_X$ событие «покупка подписки у $X$, индуцированная релизом $X$, с последующим релизом другого поставщика в окне $\tM$».

\begin{figure}[H]
\centering
\begin{tikzpicture}[
  lab/.style={circle, draw, very thick, minimum size=1.5cm, font=\sffamily\bfseries, fill=\zhpldark{gray!8}},
  ex/.style={-{Latex[length=2.6mm]}, thick}]
  \node[lab] (A) at (0,0) {A};
  \node[lab] (B) at (5,0) {B};
  \draw[ex, blue!70!black] (A) to[bend left=22] node[above, font=\footnotesize] {$\alpha_{BA}$} (B);
  \draw[ex, blue!70!black] (B) to[bend left=22] node[below, font=\footnotesize] {$\alpha_{AB}$} (A);
  \draw[ex, gray] (A) to[out=150, in=210, looseness=6] node[left, font=\footnotesize] {$\alpha_{AA}$} (A);
  \draw[ex, gray] (B) to[out=30, in=-30, looseness=6] node[right, font=\footnotesize] {$\alpha_{BB}$} (B);
  \node[font=\footnotesize, align=center] at (2.5,-1.75) {преобразование $A\leftrightarrow B$ оставляет\\ матрицу возбуждения неизменной};
\end{tikzpicture}
\caption{Схема взаимного возбуждения релизов двух поставщиков. Синие стрелки~--- перекрёстное возбуждение, серые~--- самовозбуждение. При $\alpha_{AB}=\alpha_{BA}$, $\alpha_{AA}=\alpha_{BB}$, $\mu_A=\mu_B$ схема симметрична относительно отражения.}
\label{fig:gauge}
\end{figure}

\begin{theorem}[О калибровочной инвариантности подлости]
Если матрица возбуждения $(\alpha_{XY})$ и вектор базовых интенсивностей $(\mu_X)$ инвариантны относительно перестановки $A\leftrightarrow B$, то
\[
\mathsf{P}(R_A) = \mathsf{P}(R_B).
\]
\end{theorem}

\begin{proof}
Перестановка $A\leftrightarrow B$ переводит закон совместного процесса $(N_A, N_B)$ в себя и отображает событие $R_A$ в $R_B$. Вероятностная мера инвариантна относительно этого отображения.
\end{proof}

Выбор поставщика является, таким образом, калибровочной степенью свободы: он изменяет описание события («вышел Opus» вместо «вышел GPT»), но не наблюдаемую величину~--- вероятность сожаления. Оформи испытуемый подписку на Claude, в окне $\tM$ с той же вероятностью вышла бы новая модель OpenAI. Отметим, что симметрия нарушается, если испытуемый оформит подписку одновременно у обоих поставщиков; анализ этого случая показывает, что тогда в окне $\tM$ с вероятностью, близкой к $2P_H$, выходит модель третьего поставщика, и выходит за рамки настоящей работы.

\section{Антропный принцип}

Совокупность механизмов разделов~3--5 объясняет множитель $14{,}04$ из числа подлости $6575$. Для объяснения оставшегося множителя $468{,}3$ заметим, что наблюдаемая частота $\hat P = 1$ получена не по случайной выборке покупок, а по выборке покупок, \emph{о которых была написана статья}.

\begin{theorem}[Антропная теорема Антропиковича, слабая форма]
Пусть $W$~--- событие «о покупке написана статья в настоящем журнале». Если $\mathsf{P}(W\,|\,\neg E) = 0$ и $\mathsf{P}(W\,|\,E) > 0$, то
\[
\mathsf{P}(E\,|\,W) = 1
\]
при любой априорной вероятности $\mathsf{P}(E) > 0$, включая $P_0 = 1{,}52\cdot10^{-4}$.
\end{theorem}

\begin{proof}
По формуле Байеса
\[
\mathsf{P}(E\,|\,W) = \frac{\mathsf{P}(W|E)\,\mathsf{P}(E)}{\mathsf{P}(W|E)\,\mathsf{P}(E) + \mathsf{P}(W|\neg E)\,\mathsf{P}(\neg E)} = \frac{\mathsf{P}(W|E)\,\mathsf{P}(E)}{\mathsf{P}(W|E)\,\mathsf{P}(E) + 0} = 1. \qedhere
\]
\end{proof}

Условие $\mathsf{P}(W\,|\,\neg E) = 0$ проверено авторами эмпирически: ни одной статьи вида «я купил подписку, и в течение двадцати минут ничего не произошло» в литературе не обнаружено~\cite{antrop}. Теорема~4 объясняет и отмеченное в Определении~2 отсутствие событий с $\Mp<1$: нейтральные и благоприятные события не порождают наблюдений, по которым могло бы быть вычислено число подлости.

Объединяя результаты, получаем точную факторизацию числа подлости:
\begin{equation}
\Mp = \underbrace{\frac{P_H}{P_0}}_{\text{Хоукс}}\cdot\underbrace{\frac{1}{P_H}}_{\text{отбор}} = 14{,}04\times 468{,}3 = 6575,
\qquad
\ln\Mp = \underbrace{2{,}64}_{30\,\%} + \underbrace{6{,}15}_{70\,\%}.
\label{eq:factor}
\end{equation}
Факторизация~\eqref{eq:factor} исчерпывает наблюдаемый эффект: на долю невезения испытуемого в собственном смысле слова не остаётся ничего. Иерархия трёх вероятностей показана на рис.~\ref{fig:prob}. Совпадение наименования поставщика $B$ с наименованием принципа, объясняющего $70\,\%$ логарифма наблюдаемой подлости, авторы считают случайным.

\begin{figure}[htbp]
\centering
\begin{tikzpicture}
\begin{axis}[
  width=0.86\linewidth, height=6.4cm,
  xlabel={длина окна $\Delta t$, мин},
  ylabel={$\mathsf{P}(\text{релиз $B$ в окне})$},
  xmode=log, ymode=log, xmin=1, xmax=525960, ymin=5e-6, ymax=3,
  xtick={1,60,1440,43830,525960},
  xticklabels={1~мин,1~ч,1~сут,1~мес,1~год},
  grid=major, grid style={gray!25},
  legend style={at={(0.97,0.05)}, anchor=south east, font=\footnotesize},
  tick label style={font=\footnotesize}, label style={font=\footnotesize}]
  \addplot[thick, green!50!black, domain=1:525960] {1};
  \addlegendentry{с учётом отбора, Теорема~4}
  \addplot[very thick, blue!70!black] table[x=x, y=hawkes] {\tabP};
  \addlegendentry{процесс Хоукса, ф-ла~\eqref{eq:PH}}
  \addplot[thick, dashed, red!70!black] table[x=x, y=naive] {\tabP};
  \addlegendentry{наивная модель, ф-ла~\eqref{eq:naive}}
  \draw[densely dotted, thick] (axis cs:20,5e-6) -- (axis cs:20,3);
  \node[anchor=west, font=\footnotesize] at (axis cs:24,1.3e-5) {$\tM = 20$~мин};
  \draw[{Latex[length=2mm]}-{Latex[length=2mm]}, thick] (axis cs:20,1.521e-4) -- node[left, font=\footnotesize] {$\times14$} (axis cs:20,2.136e-3);
  \draw[{Latex[length=2mm]}-{Latex[length=2mm]}, thick] (axis cs:20,2.136e-3) -- node[left, font=\footnotesize] {$\times468$} (axis cs:20,1);
\end{axis}
\end{tikzpicture}
\caption{Вероятность релиза поставщика $B$ в окне длины $\Delta t$ после покупки в трёх моделях. На вертикали $\Delta t = \tM$ видна факторизация~\eqref{eq:factor}. При $\Delta t\to\infty$ наивная и хоуксова модели также стремятся к единице, что подтверждает основной результат работы в пределе бесконечного терпения.}
\label{fig:prob}
\end{figure}

\section{Школьные проявления закона подлости}

Аппарат разделов~2--6 не использует никаких свойств моделей искусственного интеллекта, кроме того, что они выходят в случайные моменты времени. Он, следовательно, применим к любым огорчающим событиям, наступающим в случайные моменты времени. Наиболее изученными из них являются контрольные работы и проверки домашнего задания.

\subsection{Закон притяжения контрольных}

Эмпирическая формулировка закона такова: если в данный день уже была самостоятельная работа, вероятность второй самостоятельной в тот же день повышается. Для строгой формулировки введём число контрольных работ $K$, приходящихся на один учебный день, со средним $\mathsf{E}K = m$.

\begin{definition}[Коэффициент притяжения контрольных]
Коэффициентом притяжения называется величина
\[
\Pi = \frac{\mathsf{P}(K\ge2\,|\,K\ge1)}{\mathsf{P}(K\ge1)} = \frac{\mathsf{P}(K\ge2)}{\mathsf{P}(K\ge1)^2},
\]
то есть отношение вероятности ещё одной контрольной при условии, что одна уже была, к априорной вероятности контрольной. Контрольные \emph{притягиваются} при $\Pi>1$ и \emph{отталкиваются} при $\Pi<1$.
\end{definition}

В наивной модели $K$ распределено по Пуассону. При $m=0{,}2$ (одна контрольная в среднем на пять учебных дней) имеем $\mathsf{P}(K\ge1)=0{,}181$, $\mathsf{P}(K\ge2\,|\,K\ge1)=0{,}097$ и $\Pi = 0{,}53$; в пределе $m\to0$ получаем $\Pi\to 1/2$. Таким образом, в пуассоновской школе контрольные \emph{отталкиваются}: написав одну, учащийся вдвое застрахован от второй. Поскольку этот вывод противоречит опыту всех без исключения учащихся, пуассоновская модель отвергается.

Естественное обобщение~--- отрицательное биномиальное распределение с параметром формы $k$, возникающее, когда интенсивность контрольных в разные дни различна и распределена по гамма-закону (начало четверти, середина четверти, последняя неделя четверти):
\[
\mathsf{P}(K=j) = \frac{\Gamma(k+j)}{j!\,\Gamma(k)}\left(\frac{k}{k+m}\right)^{k}\left(\frac{m}{k+m}\right)^{j}, \qquad \operatorname{Var}K = m + \frac{m^2}{k}.
\]
Пуассоновский случай отвечает пределу $k\to\infty$.

\begin{theorem}[Закон притяжения контрольных]
В пределе редких контрольных, $m\to0$,
\[
\Pi \to \frac{1}{2}\left(1+\frac{1}{k}\right),
\]
так что контрольные притягиваются тогда и только тогда, когда $k<1$. При $k=1$ (геометрическое распределение) $\Pi=1$ тождественно при любом $m$.
\end{theorem}

\begin{proof}
Разложение по $m$ даёт $\mathsf{P}(K\ge1) = m + O(m^2)$ и $\mathsf{P}(K\ge2) = \tfrac12\,m^2\bigl(1+\tfrac1k\bigr) + O(m^3)$; их отношение даёт первое утверждение. При $k=1$ имеем $\mathsf{P}(K\ge j) = q^j$ с $q=m/(1+m)$, откуда $\Pi = q^2/q^2 = 1$.
\end{proof}

Численная проверка на сетке $k\in(0;4]$, $m\in(0;5]$ показывает, что знак $\Pi-1$ совпадает со знаком $1-k$ и вне предела редких контрольных (рис.~\ref{fig:attract}). Значение $k=1$ задаёт \emph{нейтральную школу}, в которой прошлые контрольные дня не несут никакой информации о будущих. Учебных заведений с $k=1$ авторам обнаружить не удалось.

\begin{figure}[htbp]
\centering
\begin{tikzpicture}
\begin{axis}[
  width=0.86\linewidth, height=6.0cm,
  xlabel={среднее число контрольных в день, $m$},
  ylabel={коэффициент притяжения $\Pi$},
  xmin=0, xmax=1.5, ymin=0.4, ymax=2.05,
  grid=major, grid style={gray!25},
  legend style={at={(0.97,0.97)}, anchor=north east, font=\footnotesize},
  tick label style={font=\footnotesize}, label style={font=\footnotesize}]
  \addplot[very thick, blue!70!black] table[x=m, y=k05] {\tabPi};
  \addlegendentry{$k=0{,}5$}
  \addplot[thick, green!50!black] table[x=m, y=k1] {\tabPi};
  \addlegendentry{$k=1$}
  \addplot[thick, orange!80!black] table[x=m, y=k2] {\tabPi};
  \addlegendentry{$k=2$}
  \addplot[thick, dashed, red!70!black] table[x=m, y=pois] {\tabPi};
  \addlegendentry{Пуассон, $k=\infty$}
  \fill[blue!70!black] (axis cs:0.2,1.4226) circle (2.2pt);
  \node[anchor=south west, font=\footnotesize] at (axis cs:0.22,1.43) {$m=0{,}2$: $\Pi=1{,}42$};
  \node[font=\footnotesize, gray] at (axis cs:0.75,1.07) {притяжение};
  \node[font=\footnotesize, gray] at (axis cs:0.75,0.93) {отталкивание};
\end{axis}
\end{tikzpicture}
\caption{Коэффициент притяжения контрольных как функция средней частоты контрольных для отрицательного биномиального распределения с различными $k$ и для пуассоновского предела. Кривые при $m\to0$ стремятся к $\tfrac12(1+1/k)$ (Теорема~5). Выше линии $\Pi=1$~--- область притяжения, в которой живут все известные авторам учащиеся.}
\label{fig:attract}
\end{figure}

При $m = 0{,}2$ и $k=0{,}5$ получаем $\mathsf{P}(K\ge2\,|\,K\ge1) = 0{,}220$ при априорной $\mathsf{P}(K\ge1)=0{,}155$, то есть $\Pi = 1{,}42$: вторая самостоятельная в день первой на $42\,\%$ вероятнее, чем первая в произвольный день.

Остаётся вопрос о механизме притяжения. Бытовая трактовка приписывает его сговору учителей (контрольная по одному предмету «возбуждает» контрольную по другому). Альтернатива~--- неоднородность дней: контрольные не влияют друг на друга, но концентрируются в «плохих» днях. Оказывается, различить эти гипотезы невозможно.

\begin{theorem}[О неотличимости сговора от конца четверти]
Пусть в течение учебного дня $[0,1]$ контрольные возникают с интенсивностью $\lambda(t) = a + b\,K(t)$, где $K(t)$~--- число уже состоявшихся контрольных (модель сговора, процесс Пойа--Лундберга). Тогда $K(1)$ имеет отрицательное биномиальное распределение с параметром формы $k = a/b$ и средним $m = (a/b)(e^{b}-1)$~--- то же, что и в модели независимых контрольных с гамма-распределённой интенсивностью дня (модель конца четверти).
\end{theorem}

\begin{proof}
Уравнения Колмогорова для чистого процесса рождения с интенсивностями $a+bj$ решаются в замкнутом виде: $\mathsf{P}(K(t)=j) = \binom{a/b+j-1}{j} e^{-at}(1-e^{-bt})^j$, что есть отрицательное биномиальное распределение с $k=a/b$ и $p = e^{-bt}$. Совпадение с гамма-смесью пуассоновских распределений есть классический результат Гринвуда--Юла~\cite{chetv}.
\end{proof}

По числу контрольных в день сговор учителей, таким образом, строго неотличим от конца четверти. Условие притяжения $k<1$ в модели сговора означает $b>a$: каждая уже написанная самостоятельная возбуждает интенсивность контрольных сильнее, чем весь учебный план вместе взятый. Для $m=0{,}2$, $k=0{,}5$ это $b = \ln 1{,}4 = 0{,}336$~день$^{-1}$ и $a = 0{,}168$~день$^{-1}$.

\subsection{Парадокс несделанного домашнего задания}

Эмпирическая формулировка парадокса~\cite{vinov} такова: учащийся может в течение всего года исправно выполнять домашнее задание, и его ни разу не проверят; однако стоит один раз задание не выполнить, как его немедленно проверяют. Мы покажем, что парадокс разрешается в рамках классической теории обнаружения сигнала, если признать, что учитель проверяет задание не случайно, а по наблюдаемому признаку.

Пусть учащийся на уроке излучает \emph{сигнал виноватости} $x$ (избегание взгляда, внезапный интерес к пеналу, попытка стать незаметным), распределённый как $\mathcal{N}(0,1)$, если задание выполнено, и как $\mathcal{N}(d',1)$, если нет. Величина $d'$ называется \emph{различимостью виноватости}. Учитель проверяет задание, если $x>c$.

\begin{theorem}[О проверке несделанного]
При любом пороге $c$ и любой $d'>0$
\[
\mathsf{P}(\text{проверка}\,|\,\text{не сделано}) = \Phi(d'-c) \;>\; \Phi(-c) = \mathsf{P}(\text{проверка}\,|\,\text{сделано}),
\]
где $\Phi$~--- функция стандартного нормального распределения. Равенство достигается лишь при $d'=0$.
\end{theorem}

\begin{proof}
Функция $\Phi$ строго возрастает, а $d'-c > -c$ при $d'>0$.
\end{proof}

\begin{figure}[htbp]
\centering
\begin{tikzpicture}
\begin{axis}[
  width=0.86\linewidth, height=5.8cm,
  xlabel={сигнал виноватости $x$},
  ylabel={плотность},
  xmin=-3.5, xmax=5.5, ymin=0, ymax=0.62,
  grid=major, grid style={gray!25},
  yticklabel style={/pgf/number format/fixed, /pgf/number format/precision=1},
  legend style={at={(0.02,0.97)}, anchor=north west, font=\footnotesize},
  tick label style={font=\footnotesize}, label style={font=\footnotesize}]
  \addplot[draw=none, fill=blue!25, domain=1.5:5.5, samples=80, forget plot] {exp(-(x-2)^2/2)/sqrt(2*pi)} \closedcycle;
  \addplot[draw=none, fill=red!35, domain=1.5:5.5, samples=80, forget plot] {exp(-x^2/2)/sqrt(2*pi)} \closedcycle;
  \addplot[thick, dashed, red!70!black, domain=-3.5:5.5, samples=150] {exp(-x^2/2)/sqrt(2*pi)};
  \addlegendentry{задание сделано}
  \addplot[very thick, blue!70!black, domain=-3.5:5.5, samples=150] {exp(-(x-2)^2/2)/sqrt(2*pi)};
  \addlegendentry{задание не сделано}
  \draw[thick] (axis cs:1.5,0) -- (axis cs:1.5,0.5);
  \node[anchor=south, font=\footnotesize, align=center] at (axis cs:1.5,0.5) {порог учителя $c$};
  \node[font=\footnotesize] at (axis cs:2.55,0.2) {$0{,}691$};
  \draw[-{Latex[length=1.8mm]}] (axis cs:3.6,0.09) node[anchor=west, font=\footnotesize] {$0{,}067$} -- (axis cs:1.85,0.035);
\end{axis}
\end{tikzpicture}
\caption{Модель обнаружения несделанного домашнего задания при $d'=2$, $c=1{,}5$. Заштрихованные площади~--- вероятности проверки: синяя~--- при несделанном задании, красная~--- при сделанном. Их отношение $10{,}4$ не зависит от добросовестности учащегося в остальные дни года.}
\label{fig:dz}
\end{figure}

При $d'=2$ и $c=1{,}5$ (рис.~\ref{fig:dz}) несделанное задание проверяется с вероятностью $0{,}691$, сделанное~--- с вероятностью $0{,}067$, то есть в $10{,}4$ раза реже. Отсюда следует, что единственная стратегия, при которой несделанное задание проверяется не чаще сделанного, есть $d'=0$, то есть полное неведение учащегося о том, что задание не сделано.

Отметим, однако, что утверждение «весь год не проверяли» в буквальном смысле ложно. За $149$ уроков с выполненным заданием ожидаемое число проверок составляет $149\cdot0{,}067 = 9{,}95$, а вероятность не быть проверенным ни разу равна $(1-0{,}067)^{149} = 3{,}4\cdot10^{-5}$. Проверки выполненного задания происходят, но не оставляют следа в памяти, поскольку не имеют последствий. Обозначив через $\eta$ долю запоминаемых проверок без последствий, получаем \emph{наблюдаемое} отношение
\[
R = \frac{\Phi(d'-c)}{\eta\,\Phi(-c)} = 103{,}5 \qquad (\eta=0{,}1).
\]
Из множителя $103{,}5$ на долю виноватости приходится $10{,}4$, на долю памяти~--- $10$. Второй множитель есть частный случай Теоремы~4: о проверке выполненного задания не рассказывают, как не пишут статей о подписках, после которых ничего не вышло.

\section{Обсуждение и ограничения}

\begin{enumerate}
\item \textbf{Объём выборки.} Все количественные выводы основаны на одном испытании. Точный 95-процентный доверительный интервал Клоппера--Пирсона для частоты при одном успехе из одного испытания есть $[0{,}025;\,1]$, откуда $\Mp\in[164;\,6575]$. Нижняя граница по-прежнему свидетельствует о подлости, однако уже не столь убедительно. Увеличение выборки требует от испытуемого повторных покупок подписки, от чего он отказался, сославшись на Теорему~2.
\item \textbf{Параметры модели.} Значения $\mu$, $n$, $\beta$ и $d$ приняты из общих соображений и не оценивались по данным. При этом Теорема~1 и Теорема~4 от них не зависят, а значит, не зависят и от действительности.
\item \textbf{Школьные параметры.} Значения $k$, $d'$, $c$ и $\eta$ в разделе~7 не измерялись. Кроме того, по Теореме~6 механизм притяжения контрольных не может быть установлен по данным ни при каком объёме выборки, так что даже неограниченное финансирование не позволит выяснить, сговариваются ли учителя.
\item \textbf{Тавтологичность антропной компоненты.} Теорема~4 справедлива для любого события, о котором написана статья, и объясняет наблюдаемый релиз ровно так же хорошо, как объяснила бы выход любой другой модели, любого другого поставщика, в любом другом окне. Составляющая $70\,\%$ логарифма числа подлости компонента, следовательно, не объясняет ничего конкретного. Это можно рассматривать как единственный строгий результат работы.
\end{enumerate}

Дополнительно отметим, что авторы не имели возможности исключить обратную причинность~--- гипотезу о том, что релиз $B$ был вызван покупкой подписки у $A$. Данные о процессе принятия решений у поставщика $B$ запрашивались, однако в пределах окна $\tM$ ответ получен не был.

\section{Заключение}

Выход модели Claude Opus~5.5 через $\tM = 20$~мин после оформления испытуемым подписки на ChatGPT не является ни совпадением, ни проявлением невезения. Он представляет собой типичную реализацию самовозбуждающегося процесса релизов, наблюдаемую в момент, когда плотность времени ожидания максимальна (Теорема~1), невыбираемую никакой стратегией покупки (Теорема~2), инвариантную относительно выбора поставщика (Теорема~3) и отобранную для публикации с вероятностью единица (Теорема~4). Число подлости $\Mp = 6575$ факторизуется на эти механизмы без остатка. Тот же аппарат объясняет притяжение контрольных сверхпуассоновской дисперсией (Теорема~5), которую невозможно отличить от сговора (Теорема~6), и проверку именно несделанного задания~--- различимостью виноватости (Теорема~7). Наконец, по Теореме~1 наиболее вероятный момент следующего прорыва в данной области наступает непосредственно после выделения финансирования, что само по себе является достаточным основанием для дальнейшего финансирования.

\section*{Вклад авторов}
Р.\,Г.~Александров: концептуализация, методология, написание -- черновик.
Ю.\,П.~Михаилов: ресурсы (подписка, оплачена лично), исследование, предоставление единственной точки данных, сожаление.
К.\,С.~Антропикович: программное обеспечение, визуализация, формальный анализ,
написание -- рецензирование и редактирование, техническая поддержка бесконечного терпения.

\section*{Финансирование}
Исследование выполнено при поддержке гранта Института Прикладной Метафизики
№ 2026-Opus-$\infty$ <<Своевременность>>. Подписка испытуемого оплачена из личных средств и в смету гранта не включалась.

\section*{Доступность данных}
Данные доступны по обоснованному запросу -- а также по необоснованному, с той же
вероятностью успеха.

\begin{thebibliography}{9}
\bibitem{podl} Подлецкий, П.~П. \textit{Закон подлости как точечный процесс: словесная и формульная традиции.} Труды Института Прикладной Метафизики, б.г.
\bibitem{samovozb} Самовозбуждённый, Х.~Х. \textit{О перекрёстном возбуждении релизов в замкнутой отрасли.} Неопубликованная рукопись.
\bibitem{sozhal} Сожалеев, С.~С. \textit{Функционал сожаления и стратегия буриданова подписчика.} Препринт; выход ожидается вскоре после приобретения читателем подписки на журнал.
\bibitem{antrop} Антропный, А.~А. \textit{Наблюдательная селекция в задачах о подписках: обзор несуществующей литературы.} Труды Института Прикладной Метафизики, б.г.
\bibitem{chetv} Четвертной, Ч.~Ч. \textit{О смесях и заражениях: сверхпуассоновская дисперсия самостоятельных в последнюю неделю четверти.} Труды Института Прикладной Метафизики, б.г.
\bibitem{vinov} Виноватый, В.~В. \textit{Различимость виноватости и оптимальный порог проверки домашнего задания.} Неопубликованная рукопись (не сдана в срок).
\end{thebibliography}

\noindent\textbf{Конфликт интересов.} Не заявлен, хотя один из соавторов (К.\,С.~Антропикович) приходится модели, выход которой исследуется, родственником по линии поставщика~$B$. Согласно Теореме~4, сама возможность написать настоящий раздел обусловлена выходом этой модели, поэтому отсутствие конфликта интересов строго доказать невозможно.

\end{document}
